\section{The original common column selection and its linked Bundle}
\label{sec:common-selected}

\status{CSR1--CSR22 reconstruct the finite column constructor
\texttt{SelC.cols\_of\_witnesses}, the exact placement and several
geometric transports, and the same-record \texttt{BWit}/\texttt{BunU}
installation. The original point, readings, multiplicity spaces and
geometric witnesses remain fixed. The final installation is conditional
on the expressly stated original division and eigensystem producers.
It does not construct the common point, prove all the published
representation-theoretic hypotheses, or establish the separate
\texttt{ratG}/\texttt{MLat} branch.}

\paragraph{Original objects and choice order.}
Use the source's true column datum \(I:\mathrm{TC}(B,X)\), with
\(I.\mathrm{Ok}\), its arithmetic \(\mathrm{Pins}\), laws \(l\), later
stages \(Z\), canonical-index condition \(I.\mathrm{IsTw}\), and
\(\mathrm{JetsIn}(I)\). Write \(D=I.\mathrm{rec7}(l,Z)\) and
\(K=B.\mathrm{LL}\). The fields, involution, primes, embeddings, words,
weights and index matrices below are these original objects. In the
application they are projections of the supplied
\(q:\mathrm{WPt}(p.s,\mathcal K)\), with
\[
 t=q.x.\mathrm{what}_2,\qquad
 R_0=t.r.z.R,\quad F_0=t.r.z.F,\quad
 Y_0=t.r.Y,\quad Z_0=t.r.Z .
 \tag{CSR.1}
\]
\(R_0\) is a real/geometric datum, whereas \(R:\mathrm{RecT}(q.O.ex)\)
is the \emph{already supplied} final record with
\(\mathrm{RecPinE}(q,R)\). These two uses of \(R\) must not be confused.
The latter pin includes \(R.T.\mathrm{cols}=q.x.\mathrm{col}.C\), the
original theta package, character frame and tuple, and the standard
kernel. The additional \(\mathrm{RXs}(t)\) supplies the all-tuple
ambient equivalences and their exact reading/period agreements; its
existence is an upstream input. For \(p=\mathrm{prmR}(N,\mathrm{SExt.none},
\Gamma)\) it is the stored \(q.x.rx\).

We first choose a column face and one \(\mathrm{ReadC}\), before any
later \(\mathrm{TwArgs}\). In the application the spectral package
\((sP,sS)\) is the transport of the \emph{whole} \(q.\mathrm{sat}\)
package along \(R.T.\mathrm{cols}=q.x.\mathrm{col}.C\). Thus its
Hilbert space, good operator, evaluation point \(q.\mathrm{sat}.t\),
positive scalar \(q.\mathrm{sat}.c\) and periods are not chosen afresh.
The notation \(A=B\) between dependent records means equality after
the displayed equality of their indexing data; \texttt{HEq} in the
source records precisely this transport.

% BEGIN CSR1
\paragraph{CSR1: one literal level-row reader.}\label{claim:CSR1}
Let \(\sigma:K\hookrightarrow\mathbb C\), let
\(j:K^d\to V\) be \(\sigma\)-semilinear, and suppose the images
\(j(e_a)\) are complex-linearly independent and span \(V\).
For every specified \(K\)-linear equivalence
\(\Psi:K^d\simeq K^h\), there is a unique complex-linear equivalence
\(\iota:V\simeq\mathbb C^h\) satisfying
\[
                  \iota(j(y))=\sigma(\Psi y).
 \tag{CSR.2}
\]
Indeed, use \(j(e_a)\) as a complex basis. If \(A\) is the matrix of
\(\Psi\), define \(\iota\) by \(\sigma(A)\) on its coordinates.
The inverse matrix is \(\sigma(A^{-1})\). This proves both
invertibility and (CSR.2), including the zero-dimensional case.
If \(j\) originally takes values in \(A_a\), an equality \(a=b\)
and a specified \(E:V\simeq A_b\) replace it by
\(E^{-1}\circ j\); no new coordinate choice is involved.

Apply this to the original row level. Write \(e_P\) for its finite
character-averaging idempotent. The original matrices
\(\mathrm{lvlE}\) and \(\mathrm{lvlPr}\) express a basis of
\(\operatorname{im}(e_P)\) and its coordinate projection; in particular
\(\mathrm{lvlPr}\,\mathrm{lvlE}=1\). Restriction of an equivariant
function to the fixed class representatives is injective. Its image
at the prescribed level is exactly \(\operatorname{im}(e_P)\):
averaging fixes a function of that level, and an averaged function
has the required transformation law. Thus
\[
 y\longmapsto
 \mathrm{basis}^{-1}(\mathrm{lvlE}\,y)
 \tag{CSR.3}
\]
is an equivalence from \(K^{\mathrm{lvlDim}}\) to the \emph{actual}
level row functions. For the base tuple its twist-prime set is empty,
so the original row cast identifies this space with
\(\mathrm{rowGSD}\); composing (CSR.3), this cast and its stored
basis coordinates is the specified \(\Psi\) in (CSR.2).
The finiteness and form laws used here are TL2, TL6--TL7 and TG3
on the supplied \(I\); no new finiteness assertion is assumed.

For the family version, let \(\mathrm{tup}\) be the original injective
tuple map, with \(\mathrm{tup}(\mu_0)=b\), and take the supplied
ambient equivalence \(E\) at \(b\). If \(\mathrm{tup}(t)=b\), then
injectivity gives \(t=\mu_0\). Hence the \emph{one} reader just
constructed works for every such \(t\); a choice varying with a
proof of \(\mathrm{tup}(t)=b\) is unnecessary.
It has both clauses of \texttt{exists\_read\_fam}. For the second,
put \(y=\mathrm{lvlPr}\,\mathrm{vE}\,F\). The restriction in
(CSR.3) becomes the restriction of the source's \(\mathrm{evC}(F)\),
because \(\mathrm{vE}F\) already lies in the level part and
\(\mathrm{lvlE}\,\mathrm{lvlPr}\) is the identity there.
Injectivity of restriction therefore identifies the entire function
with \(\mathrm{evC}(F)\), not just its class coordinates.
This proves the exact two evaluation formulas required later by
\texttt{columnEvaluationOk}.
% END CSR1

% BEGIN CSR2
\paragraph{CSR2: character membership and the original signed weights.}
\label{claim:CSR2}
For a finite group and an irreducible characteristic-zero character
\(\chi\), the central projector belonging to a finite set \(X_g\)
of irreducible characters has trace
\[
 \operatorname{tr}\rho_\chi(e_{X_g})=
 \begin{cases}\chi(1),&\chi\in X_g,\\0,&\chi\notin X_g.\end{cases}
 \tag{CSR.4}
\]
To see this, expand the projector as the sum of the usual
\(\dim(\pi)/|G|\)-weighted character sums. Orthogonality follows
by averaging a linear map between two irreducible modules: its
average is an intertwiner, is zero between distinct modules, and
is scalar on an irreducible module; its trace determines that scalar.
Apply this to matrix units and sum the diagonal identities.
Because \(\chi(1)\) is a positive integer in characteristic zero,
nonzero trace is equivalent to membership. Reindexing this finite
sum through the \emph{specified} group isomorphism and applying the
specified coefficient embedding proves the source's
\texttt{mem\_iff\_XgE}. Its hypothesis about the central idempotent
is exactly the \(e_0\) clause of the supplied division group tie.

For weights, put \(\lambda_{u,i,\tau}=B.\mathrm{wts.lamTrue}(u,i,\tau)\)
and \(s_{u,\mu,i,\tau}=(D.\mathrm{chTw}\,u\,\mu).\mathrm{sSharp}(i,\tau)\).
The original weight definitions and signed character identities give,
for \(i<m\) and every old coordinate \(r<N_i\),
\[
 \begin{split}
 \nu_T(u,i,\tau,r)&=\lambda_{u,i,\tau}(r)+s_{u,\mu,i,\tau},\\
 \nu_{\rm last}(u,i,\tau)
   &=\operatorname{detShift}(\lambda_{u,i,\tau})+s_{u,\mu,i,\tau},\\
 \operatorname{detShift}(\lambda_{u,i+1,\tau})
   &=(D.\mathrm{chTw}\,u\,\mu).\mathrm{sDiv}(i,\tau).
 \end{split}                                                    \tag{CSR.5}
\]
Here the first equality is an equality of integers, obtained from the
source's equality of rational/real weights by injectivity of those
embeddings. The last coordinate exists since \(N_i>0\), and
\(N_i<N_{i+1}\), so at step \(i+1\) it belongs to the newly added
block. Its defining constant weight is precisely \(\mathrm{sDiv}\).
Dominance gives
\(\lambda(r)-\operatorname{detShift}(\lambda)\ge0\). Consequently
\[
 (\lambda(r)-\operatorname{detShift}\lambda)^{+}+\nu_{\rm last}
       =\nu_T(r),\qquad a^{+}=\max(a,0)\in\mathbb N .
 \tag{CSR.6}
\]
The fourth \(\mathrm{WtTie}\) clause is the original
jet's weight-space membership: it is the supplied
\(\mathrm{JetsIn}\) assertion with the unchanged canonical chain.
Thus (CSR.5)--(CSR.6), together with that supplied clause, give all
four fields of \texttt{wtTie\_base}. The seam in
\(\mathrm{StdKer}\) is the original \(\mathrm{seamOf}\); on the
original embedding index its \(\mathrm{iT}^{-1}\) is the identity.
There is no additional permutation or change of sign in these
formulas. In particular the resulting \(\mathrm{lamOfCol}\) is
literally \(B.\mathrm{wts.lamTrue}\).
% END CSR2

% BEGIN CSR3
\paragraph{CSR3: the actual ramified coordinate equivalence and metric.}
\label{claim:CSR3}
Fix a word \(u\). Let \(\mathcal C_u^{\rm raw}\) be the original
rational column-function space and let \(\mathcal C_u^{0}\)
be the actual zero-step form space \((I.\mathrm{fm}\,u).\mathrm{lat0}.V\).
Let \(\kappa:\mathcal C_u^{\rm raw}\to\mathcal C_u^{0}\)
send a column function to its value function on the source's
preferred zero-step decorations.
At step zero the index is unique and the coupling set is the fixed
empty one. Forgetting the extra frames gives a bijection of classes:
lift a class by choosing its frames; two lifts of the same class
differ by the allowed right-frame transformations. Equivariance then
makes \(\kappa\) and the inverse extension well-defined.
The original integral value conditions coincide, and the
stabilizer groups coincide. Thus the finite class sums defining
every archimedean norm coincide term by term. This is the zero-step
specialization of TL4--TL7 and TX10, with exactly their preferred
decorations, not an abstract unitary equivalence.

The diamond operators act on the same residue coordinates under
\(\kappa\). Therefore their joint norm kernels match:
\[
 \kappa(\mathrm{MramV}_u)=\mathrm{ramSub}_u .
 \tag{CSR.7}
\]
The original later stage \(Z.V\) supplies an injective map
\(\nu\mathrm{In}_u:(Z.V.\mathrm{div.Nu}\,u).V\to\mathcal C_u^{0}\)
with image exactly \(\mathrm{ramSub}_u\). Its lattice and norm
clauses are part of that \emph{supplied} stage, as in SE5.
Define the map \(e_M\) by
\[
                  \kappa(e_M(n))=\nu\mathrm{In}_u(n).
 \tag{CSR.8}
\]
The injectivity and range statements give a unique value and an
inverse on every ramified column. Additivity and \(K\)-linearity
follow by applying the injective \(\kappa\). Equality of the actual
integral lattices follows from their unchanged pointwise value
conditions.

Let \(b_u=\mathrm{ramBas}_u\), the original basis of
\(\mathrm{MramV}_u\), and define
\[
 e_u:K^{\mathrm{ramDim}_u}\simeq(Z.V.\mathrm{div.Nu}\,u).V,
 \qquad e_u(x)=e_M^{-1}(b_u x).
 \tag{CSR.9}
\]
Then
\(\nu\mathrm{In}_u(e_u x)=\mathrm{colForm}(u,\mu_0,
\mathrm{colOfCoord}(x))\) exactly. Its norm at every infinite place
equals the norm of the original \(\mathrm{MramH}\) coordinate vector:
use the supplied \(\nu\mathrm{In}\) norm law, the class-sum equality
above, and the fact that \(b_u\) is the basis used in
\(\mathrm{MramH}\).

For clarity, equality of these diagonal norms gives the \emph{whole}
Hermitian form. For any \(x,y\), use \(x+y\) to determine the real
part of their pairing. At a chosen infinite place \(v\), choose
\(a\in K\) with \(\operatorname{Im}v(a)\ne0\), possible because the
original coefficient field is totally complex. The norm of \(x+ay\),
together with the already known real part, determines its imaginary
part. This proves
\[
 \operatorname{majAt}_{\mathrm{Nu},v}(e_u x,e_u y)
      =\operatorname{majAt}_{\mathrm{MramH},v}(x,y)
 \quad\text{for all }v,x,y.                                    \tag{CSR.10}
\]
There is no assumption that an arbitrary embedding commutes with
conjugation: the two actual norms are compared first, and
polarization is carried out separately at each place.
% END CSR3

% BEGIN CSR4
\paragraph{CSR4: unchanged kernel entries and the literal period point.}
\label{claim:CSR4}
Let \(k\) be the supplied kernel with \(\mathrm{KPin}(q,k)\).
We first derive its character equality from its literal relative
determinant pin. If \(\mathfrak a\bar{\mathfrak a}=\mathcal O_E\),
then the rank-one lattice \(\mathfrak a\subset E\) is full,
finitely generated and self-dual. Invertibility of the
fractional ideal gives finite generation. It contains a
nonzero element, hence spans \(E\); and its Hermitian dual is
\(\bar{\mathfrak a}^{-1}=\mathfrak a\).
The determinant ideal relative to \(\mathcal O_E\) is
\(\mathfrak a\) itself, since its determinants are the
one-by-one entries.

Apply the final \(\mathrm{RelDetPinned}\) field to
\(\mathcal O_E\) and \(\mathfrak a\). It identifies the
\(\mathrm{idl}\)-image of their actual relative determinant
class with \(\mathfrak a\). Injectivity of \(\mathrm{idl}\)
therefore forces the preimage used by \(\mathrm{chiRec}\)
to be this same class. The finite-character formula of
HF14--HF15 evaluates it as the product of the original
local character values with exponents
\(\max(\operatorname{ord}_P\mathfrak a,0)\); that is
exactly \(I.\chi(\mathfrak a)\). The finer source quotient
has the same character values because its additional
relations lie in the character kernel, also checked in
HF14--HF15. Finally the kernel's \(\chi_{\rm rec}\)
field gives \(k.\chi(\mathfrak a)=I.\chi(\mathfrak a)\).
Thus the character equality has been derived, not added
as a new \(\mathrm{KPin}\) premise.

If \(N=N^\#\) is an original self-dual lattice in the
standard Hermitian space, its determinant ideal
\(\mathfrak a\) is norm one. To see this locally, use
the dual pairing between the lattice at \(P\) and the
conjugate lattice at \(\bar P\). Self-duality identifies
their basis matrices as inverse-transpose duals, up to
integral invertible matrices. Taking determinants gives
\(\operatorname{ord}_{P}\mathfrak a+
\operatorname{ord}_{\bar P}\mathfrak a=0\).
This is the paired-place determinant argument of HF15;
it also treats conjugation-fixed primes. Hence
\(\mathfrak a\bar{\mathfrak a}=\mathcal O_E\).
The two character factors used by the theta jet agree
for both lattices entering a kernel.

The original jet formulas then have the same integer exponent
\(e_{\rm Nat}\), the same theta-index function and the same
differentiated coefficients. Every actual kernel entry agrees.
The intermediate matrix kernels are finite families of these
entries; the two end kernels have the same conditional support
tests. Consequently the complete kernel chain is unchanged,
including its zero branches. One may therefore use the original
canonical coupling index, with its actual \(\mathrm{YbHyp}\),
\(\mathrm{ChIdx.Ok}\), form laws and evaluation datum. The latter
are supplied by \(I.\mathrm{Ok}\), \(\mathrm{IsTw}\), its pins and
the retained jets; CI17--CI18 and TL12--TL13 identify the original
canonical choices.

For a ramified coordinate vector \(x\), a used tuple \(a\), and the
actual row function \(z\), the finite coupling period is computed
by the successive kernel matrices, ending with the original
weighted row class sum. TL12 and FC8 identify this finite
composition with the period in the later stage; TX10 and TX14
give the same identification at a nonbase twist. More explicitly,
at the base tuple substitute the identity
\(\nu\mathrm{In}_u(e_u x)=\mathrm{colForm}(u,\mu_0,\phi)\)
from CSR3 into the stored base-period identity. At a nonbase
tuple apply its stored period-transport identity first and then
the same zero-step form \(\phi\). The supplied \(\mathrm{RXs}\)
equivalences identify its row argument with \(z\) pointwise.
All class weights and all kernel factors are unchanged, so
\[
 \operatorname{emb}_{K}\bigl(\mathrm{face.Pt}(u,a,x,f)\bigr)
  =D.\mathrm{perCol}\bigl(u,a,D.jN_u(e_u x),e_a(\mathrm{toCx}(z))\bigr)
 \tag{CSR.11}
\]
when \(z\) is the actual \(\mathrm{evC}(f)\).
This proves \texttt{face\_pt} with the original coupling and
periods. It does not derive the analytic jet algebraicity or the
existence of the later stages from the finite calculation.
% END CSR4

% BEGIN CSR5
\paragraph{CSR5: one original \(\mathrm{ReadC}\), with all six fields.}
\label{claim:CSR5}
Use the supplied \(\mathrm{RXs}(t)\). Its coordinate clause says
that \(Y_0.jA_a(rV_0f)\) equals \(e_a(\mathrm{toCx}(z))\) for a
row \(z\) whose restriction is the specified level coordinate.
The \(\mathrm{evC}\) restriction formula of CSR1 has exactly that
coordinate, so restriction injectivity forces
\(z=\mathrm{evC}(f)\). Its \(\mathrm{SameOn}\) identities identify
the tuple and \(jA\) with \(Y_0\), not merely a reader of equal
dimension.

For \emph{any} period family \(\mathrm{Sp}\) and scale function
\(\mathrm{Zsc}\) on the face in CSR4, put
\[
 F^*=\mathrm{reFace}(F_0,D_c,\mathrm{Sp},\mathrm{Zsc}).
\]
Define \(\mathrm{ReadC}(R_0,F^*,q.O.ex.p0,Y_0)\) as follows:
\[
 \begin{array}{c|l}
 \mathrm{wU}&u\mapsto u\\
 \mathrm{bij}&\text{the inverse is the same identity map}\\
 \mathrm{eNu}_u&e_u\text{ from (CSR.9)}\\
 \mathrm{NuC}_u&e_u^{-1}((D.\mathrm{div.Nu}\,u).L)\\
 \mathrm{nu}_u&\text{the lattice equality and all-place pairing (CSR.10)}\\
 \mathrm{pt}&\text{equation (CSR.11), with }e_a(\mathrm{toCx}z)
                  =Y_0.jA_a(rV_0f).
 \end{array}                                                    \tag{CSR.12}
\]
Surjectivity of \(e_u\) gives
\(e_u(e_u^{-1}L)=L\), so the lattice clause is equality.
The \(\mathrm{pt}\) clause holds for every word, used tuple and both
coordinate vectors. These formulas do not inspect \(\mathrm{Sp}\)
or \(\mathrm{Zsc}\). Thus one family \(e_u\), chosen once from the
original ramified bases, works for any later choice of those
period/scale fields. It remains a \(\mathrm{ReadC}\) on the
original \(MMap\), \(Y_0.jA\), \(D.jN\) and \(D.\mathrm{perCol}\).
% END CSR5

% BEGIN CSR6
\paragraph{CSR6: construct the original compatible column record.}
\label{claim:CSR6}
For the fixed kernel \(k\), define \(c=\mathrm{colRecK}(k)\).
This is the source's \(\mathrm{colRecC}\) with character \(k.\chi\),
not a replacement column family. Its data are
\[
 \begin{gathered}
 \mathrm{Pr}=B.\mathrm{primes},\quad
 \mathrm{Rk}=B.R,\quad W=B.\mathrm{wts},\quad
 \mathrm{Ch}=I.\mathrm{Ch},\quad \mathrm{sel}=I.\mathrm{sel},\\
 S_g=I.S,\quad \mathrm{Lv}=I.P_s,\quad X_g=X_g^{\,0},
 \quad m_2=m-2,\quad G_c=I.\mathrm{glob},\\
 \mu(t)=t,\quad \mathrm{cst}(t)=\mathrm{cstT}(t),\quad
 L(t)=\mathrm{idxT}(t),\quad
 \mathrm{Mem}=B.\mathrm{Word},\quad \mathrm{wd}=\mathrm{id}.
 \end{gathered}                                                   \tag{CSR.13}
\]
Its coefficient field is \(B.\mathrm{fields}\), its embedding is the
original one, and its ambient jets are \(I.J.\mathrm{amb}(B,k.\chi)\).
For each tuple the selected final coupling index is the original
canonical one. CSR4 proves that its kernel/evaluation conditions
are the original ones. The remaining proof fields are the given
field law, index law, row form law and the \(\mathrm{evC}\)
identification. This specifies every data field and the proofs
\(\mathrm{hX,hL,hP,hY,hK,hF,hE}\) of the column record.

Now let \(T:\mathrm{TiedDataSK}(ex)\) have this standard kernel,
and retain the source's actual extended-data remainder on this
family, with tuple map \(\mathrm{ixOf}\) and word map the identity.
Keep its \(Z_s,S_0,S_1,C,\mathrm{NuC}\), its \(I,JA,JM\), and all its
division, ambient and lift ties. Define a compatible column
record with the above \(c\), the given entire period family
\(\mathrm{Sp}\), \(X=B.\mathrm{fields}\), \(e=B.\mathrm{tauLL}\),
\(v_0=\mathrm{InfinitePlace.mk}(B.\mathrm{embLL})\),
\(TU=\mathrm{tuOf}(I,U)\), \(\xi=\mathrm{xiValOf}(I)\), and the
original inert transfer \(\mathrm{tr}\).
Its \(\mathrm{colTie}\) is literal face/period equality.

All fourteen entries of \(\mathrm{kerTie}\) are checked at once:
the field, rank, primes, characters, selected primes, integer set,
tuple, local constants, index matrices, \(m_2\), global constants,
ambient data, member type and word map are those pinned by
\(\mathrm{StdKer}\). For tuples, constants and indices, the
identity \(\mathrm{tup}(\mathrm{ixOf}(t))=t\) and
\(\mathrm{IsTw}\) identify the canonical \(\mathrm{cstT}\) and
\(\mathrm{idxT}\) with the source values. The word tie is the
given identity. The level and sign ties, and
\(\mathrm{famEq,gEq,plEq,i0Eq,tEq,muTie}\), are literal equalities
of the objects in (CSR.13). Its \(\mathrm{ok1}\) is the supplied
original \(\mathrm{famC.Ok3}(1)\). This constructs the full
\(\mathrm{SquareIntegrableColumnsCompatible}(T,P)\).
% END CSR6

% BEGIN CSR7
\paragraph{CSR7: all ten \(\mathrm{ColPins}\) on those columns.}
\label{claim:CSR7}
The compatible record of CSR6 has the following complete pin
certificate. Its group tie, level/character tie and point-character
tie are identities of the original groups, \(P_0,P_1,\psi_0,\psi_1\).
Its weight law is the original prime/weight law of \(B\), and its
character law is \(I.\mathrm{Ok.ch\_ok}\). Its Harish--Chandra
tie is unchanged because the standard seam acts as the identity
on the embedding index. CSR2 supplies the determinant-shift,
nonnegative-weight and jet clauses of \(\mathrm{wtTie}\), and
(CSR.4) supplies the exact \(X_g\) membership test through the
original central idempotent.

For \(\mathrm{qNF}\), use the \emph{selected} canonical index
property proved in CI17, not an existence of an equivalent
double coset. For each original tuple, permitted \(j,b\) and
nonauxiliary place, it gives integers \(n_a\), one for each new
coordinate, with \(-j\le n_a\le0\), and a vector \(z\) such that
the indicated vector \((z,0)\) belongs to
the prescribed lattice and the image of that lattice under the
actual index matrix equals its image under the original matrix
\(\mathrm{nfMat}(n,c,q,(n_a)_a,z)\). The latter has the identity on
the old \(n\)-block, \(q^{n_a}\) in new diagonal position \(a\), and the
specified \(q^{n_0}z\) column. Reindexing by
\(\mathrm{ixOf}\) preserves the same matrix and lattice.
This is exactly the source's \(\mathrm{IdxMatNF}\), including
its image equality.

Finally the empty twist-prime set at \(\mu_0\) gives the actual
\(\mathrm{rowGSD.Ok}\). For each word apply CSR1 with the stored
\(D.eMt_u\), the actual tuple map, its injectivity and base value,
and the independent spanning images supplied by
the sharp \(\mathrm{FormsTie}\) in \(P.\mathrm{tied.div}\).
This supplies one complex \(\iota_u\) with both required
restriction and evaluation formulas. Those are the two clauses
of \(\mathrm{columnEvaluationOk}\). We have now supplied
\[
 \begin{gathered}
 \mathrm{grpTie,\ psTie,\ pointCharacterEq,\ wOk,\ chOk},\\
 \mathrm{hcTie,\ wtTie,\ xgTie,\ qNF,\ columnEvaluationOk},
 \end{gathered}
\]
on the same \(P\) and same
columns, with no residual pin field.
% END CSR7

% BEGIN CSR8
\paragraph{CSR8: simultaneous transport through the whole original package.}
\label{claim:CSR8}
Suppose \(P:\mathrm{RowSideFamily}(ex)\) has the specified
equalities with the family in CSR6: group and its group structure,
place and inert index types, tuple type, \(\mathrm{A3}\), both
levels/characters, tuple and word maps, \(Z_s\), column face and
period. First substitute the group/type equalities, then their
dependent structures, then the family, levels and \(Z_s\),
and finally the face/period equalities. Finite enumerations of
the same finite type contribute no extra datum to these maps.
After these substitutions, the remaining fields of \(P\) are
exactly an \(\mathrm{ExtendedDataRemainder}\). Keep them, and
apply CSR6--CSR7. Substituting back gives compatible columns on
the \emph{original \(P\)}, with the same whole \(c,\mathrm{Sp},W,
X_g\) and face. This proves the complete transport
\texttt{exists\_rowsKC\_colPins}; no form, lattice or period
witness is independently selected during the transport.

For \(\mathrm{TwArgs}\), its \(\mathrm{pack}\) is equality of the
dependent tuple
\[
 (G,\mathrm{grp},\mathrm{fin},\mathrm{Pl}_0,I_0,T,\mathrm{A3},
    \mathrm{CoreJ},\mathrm{alg},R_{\rm real},F,c_G).
 \tag{CSR.14}
\]
Projecting (CSR.14) identifies both levels and the \emph{functions}
\(F.D_c,F.\mathrm{Sp},F.\mathrm{Zsc}\), not just one selected
value. The scalar/error update leaves the subspace
\(Z_s=(R_{\rm real}.\mathrm{toCx3}(F)).zs\) unchanged.
Consequently these identities also hold at that exact \(Z_s\).
Together with the stored \(\mathrm{tup}\) equality and the given
word-map equality, they are all the hypotheses used in the
preceding transport. This yields \texttt{exists\_cols} for the
actual \(\mathrm{divisionPair}(x)\), retaining \(x.F.\mathrm{Sp}\).
% END CSR8

% BEGIN CSR9
\paragraph{CSR9: the original column constructor.}
\label{claim:CSR9}
Assume the original \(q,k\) satisfy \(\mathrm{KPin}(q,k)\)
and that \(\mathrm{RXs}(t)\) is inhabited. Choose one such reading.
Put \(D_c=\mathrm{colRecK}(k).\mathrm{face}\).
At every \(Z_s\), choose the initial period object
\[
       \mathrm{Sp}_0(Z_s)
           =(\mathbb C,\text{its usual Hilbert structure},
                         \text{all period maps }0,\ \{0\}),
       \qquad \mathrm{Zsc}_0=F_0.\mathrm{Zsc}.
 \tag{CSR.15}
\]
Use CSR5 to construct \(C_0:\mathrm{ReadC}\) on these data.
This is one choice before all the following quantifiers:
for every \(T:\mathrm{TiedDataSK}\) with \(T.\mathrm{ker}=k\),
every original \(a:\mathrm{TwArgs}(q.O.ex,t)\), and every \(x\)
with
\[
 a.x=x,\quad a.D_c=D_c,\quad
 x.C.\mathrm{wU}=C_0.\mathrm{wU},\quad
 x.C.\mathrm{eNu}=C_0.\mathrm{eNu},
 \tag{CSR.16}
\]
there is
\(C:\mathrm{SquareIntegrableColumnsCompatible}
(T,\mathrm{divisionPair}(x))\), with
\[
 \begin{gathered}
 C.c.\mathrm{face}=x.F.D_c,\quad C.\mathrm{Sp}=x.F.\mathrm{Sp},\\
 C.c.W=q.B.\mathrm{wts},\quad
 C.c.X_g=q.x.\mathrm{what}_2.r.z.X_g,\quad \mathrm{ColPins}(C).
 \end{gathered}
 \tag{CSR.17}
\]
Indeed the standard kernel pin and (CSR.14) permit CSR8, and
the word map in (CSR.16) is the identity fixed in (CSR.12).
The displayed \(\mathrm{eNu}\) equality is part of the literal
source interface, although this last finite \(\mathrm{ColPins}\)
step does not use it; it is essential for the subsequent spectral
period installation. The zero choice in (CSR.15) constructs no
nonzero period and discards none: (CSR.17) uses each later
\(x.F.\mathrm{Sp}\) exactly.
% END CSR9

% BEGIN CSR10
\paragraph{CSR10: the exact original evaluation placement.}
\label{claim:CSR10}
Let \(H\) be the original Hilbert model, \(e:H\to X\) its
linear model map, and \(E\subset X=\mathrm{MMap}(M)\) the
\emph{actual} selected level subspace. Retain the original
\(\mathrm{models\_HE}\) and \(\mathrm{EvalTieT}(M,\nu,\mathcal F,
E,S,G,\iota,t,c)\), with \(c>0\).
Unpacking the latter gives the selected good evaluation map into
the good Hilbert subspace. The source's
\(\mathrm{sEEofTT}\) is its composition with the fixed conjugate
Hilbert isometry. Thus it is conjugate-linear and
\[
 \|\mathrm{sEEofTT}(v)\|^2=c\,\operatorname{Re}\mathrm{ipM}(v,v),
 \quad
 \mathrm{sEEofTT}(v)\in\mathrm{goodSub}.
 \tag{CSR.18}
\]
These statements follow directly from the \emph{witness chosen in}
\(\mathrm{EvalTieT}\): the conjugate isometry preserves the norm
and carries its good subspace to the prescribed good subspace.
We do not replace this chosen map by another map satisfying the
same norm formula.

For a linear \(e:H\to X\), a subspace \(E\subset X\) and a
conjugate-linear \(\phi:E\to H'\), extend the linear map
\(e:e^{-1}(E)\to E\) to \(g:H\to E\) by extending a basis of
\(e^{-1}(E)\) to one of \(H\) and setting \(g=0\) on the
added basis vectors. This is an algebraic extension; continuity
on the complement is neither claimed nor required by the source's
semilinear-map type. For the original conjugation twist
\(\mathrm{Tw}_s\), set
\[
 \widetilde\phi(x)=
   \mathrm{Tw}_s\bigl(z\,\phi(g(\mathrm{Tw}_s^{-1}x))\bigr).
 \tag{CSR.19}
\]
It is conjugate-linear because the scalar conjugation defining
\(\mathrm{Tw}_s\) commutes with ordinary complex conjugation.
On \(e^{-1}(E)\), it is precisely the prescribed formula.

Here \(n_P=\mathrm{RealDat.nP}(F,\mathrm{true})\) is the
cardinality of the original finite subgroup \(P\), so \(n_P>0\).
Choose
\[
                  \kappa=(c\,n_P)^{-1/2}>0,\qquad z=\kappa.
 \tag{CSR.20}
\]
The model identity says
\(\operatorname{Re}\mathrm{ipM}(e(v),e(v))=\|v\|^2\).
Equations (CSR.18)--(CSR.20) therefore give, for every \(v\)
whose original model image lies in \(E\),
\[
 s_Ev=\mathrm{Tw}_s\bigl((1\cdot\kappa)\,
                   \mathrm{sEEofTT}(e(v))\bigr),\qquad
               n_P\|s_Ev\|^2=\|v\|^2 .                        \tag{CSR.21}
\]
Its image on those vectors remains in the good subspace, since
that is a vector subspace. This proves the full placement
statement on the literal original \(\mathrm{MMap}\), including
its level guard and its factor \(n_P\). The extension makes no
norm claim outside that guarded subspace.
% END CSR10

% BEGIN CSR11
\paragraph{CSR11: changing inactive scalars preserves both row metrics.}
\label{claim:CSR11}
Fix the original \(R_0,F_0\) with their \(\mathrm{Pins}\) and
\(\mathrm{Holds}\), and retain the exact support condition
\[
 E_p\ne0\ \Longrightarrow\
 \left(W_{i,p}\ne0\ \text{and}\
       \forall v\ne v_0,\ \mathrm{PartAt}(F_0,i,p,v)\right)
 \quad\text{for every }i.                                    \tag{CSR.22}
\]
This is the source's \(\mathrm{parts}\) output; it is an
upstream multiplicity/representation-theoretic input here.
Given real numbers \(c_i\), define
\[
 \lambda'_{i,t,p,v}=
 \begin{cases}
 e^{c_i},&
 (v=v_0\ \&\ W_{i,p}=0)\ \text{or}\
 (v\ne v_0\ \&\ \neg\mathrm{PartAt}(F_0,i,p,v)),\\
 \lambda_{i,t,p,v},&\text{otherwise}.
 \end{cases}                                                  \tag{CSR.23}
\]
This is exactly \(\mathrm{zeroLam}\). It is positive, equals
the old scalar at the distinguished place whenever \(W_{i,p}\ne0\),
and equals it away from that place whenever the corresponding
part is present. By (CSR.22) it equals the old scalar for
every \(i,t,v\) on every nonzero \(E_p\).

We give the finite argument needed to preserve the row metric.
The rational row subspace \(Z\) is the inverse image under its
model map \(j\) of the range of the adjoint theta map.
On \(Z\), \(\Theta j\) is injective: if \(j(z)=\Theta^*w\)
and \(\Theta j(z)=0\), then
\[
       \|j(z)\|^2=\langle\Theta^*w,j(z)\rangle
                      =\langle w,\Theta j(z)\rangle=0.
 \tag{CSR.24}
\]
The supplied independent model-basis images make \(j\) injective,
so \(z=0\). The rational Hecke operators are commuting and split,
and intertwine with the operators on the theta image, by the
original \(\mathrm{Ok3}\) and primitive Hecke fields. Their
generalized-eigenspace projectors are finite polynomials in these operators.
The stable subspace \(Z\) is preserved by these projectors:
the adjoint-closed Hecke laws preserve the adjoint range, since
\(A\Theta^*=\Theta^*B\) follows by taking the adjoint of the
paired intertwining law. Decompose a vector into its finitely
many joint generalized-eigenspace components. The intertwiner preserves
each component, and (CSR.24) prevents a component in \(Z\)
from disappearing. Consequently a rational level component
of type \(p\) maps into \(E_p\), and if \(E_p=0\) that
rational component is zero. This works at both levels \(P_0,P_1\);
their level idempotents commute with the same operators.

Write \(Q_p\) for either rational component and \(\Pi_p\) for
its actual coordinate projector. If \(Q_p=0\), then
\(\Pi_p=0\), by expressing its image in a basis of \(Q_p\).
The original weighted row Gram matrix is the finite sum
\[
        \sum_p\Bigl(\kappa_v\prod_i\lambda_{i,\mu_0,p,v}\Bigr)
                           \operatorname{restrMet}(\Pi_p,X_Q).
 \tag{CSR.25}
\]
If \(\Pi_p\ne0\), the preceding paragraph gives \(E_p\ne0\),
so every coefficient in (CSR.25) is unchanged by (CSR.23).
If \(\Pi_p=0\), its summand is zero. Thus the entire matrices
coincide. The ensuing \(\mathrm{weightFrame}\), row frame and
\(\mathrm{XE}\) are the same functions of that matrix and of
unchanged data; they are equal, not merely isometric.

The conjugate scalar identity is preserved as well. When
\(\mathrm{PartAt}(F_0,i,p,v)\) fails, its defining nonvanishing
test says that every relevant diagonal majorant is zero.
A majorant is a sum of squared absolute coordinates; if its
diagonal value at \(x\) is zero, all those coordinates of \(x\)
vanish, so its pairing of \(x\) with every \(y\) is zero.
Both sides of the scalar identity are then zero regardless of
the changed scalar. On every present part the scalar was
unchanged. This proves the complete \(\mathrm{GalD}\) transport
without a new rationality or nonvanishing assumption.
% END CSR11

% BEGIN CSR12
\paragraph{CSR12: the exact re-face and reading update.}
\label{claim:CSR12}
Let \(D_c,\mathrm{Sp},\mathrm{Zsc}\) be arbitrary on the original
family. Set \(c_i=(\mathrm{Zsc}(Z_s)).l6m(i)\), take
\(\lambda'\) from (CSR.23), keep \(\kappa=R_0.\mathrm{kap}\),
and set
\[
 \begin{split}
 R'&=(R_0.\mathrm{updS}(\kappa,\lambda')).\mathrm{withErr}(e'),\\
 F'&=\mathrm{reFace}(F_0,D_c,\mathrm{Sp},\mathrm{Zsc}),\\
 e'(w,j)&=\mathrm{errOf}\bigl((R_0.\mathrm{updS}(\kappa,\lambda'))
                 .\mathrm{toCx3}(F_0).\mathrm{data}(F'.cr,\mathrm{false}),
                    F'.\mathrm{fc}(Z_s).Z.F_0,w\bigr).
 \end{split}                                                    \tag{CSR.26}
\]
The last expression is independent of \(j\), as required by the
original \(\mathrm{ErrEq}\) convention. Substituting its definition
proves that entire error equality.

Here is a field check for \(\mathrm{Pins}\) and \(\mathrm{Holds}\).
The model, inert-type, full-image, phase, all-place and tuple
fields are unaffected by a scalar/error update. For
\(\mathrm{IsScal}\), the distinguished-place equation uses
(CSR.23) on a nonzero \(W_{i,p}\); the equation away from that
place uses its \(\mathrm{PartAt}\) branch. The two absent-part
clauses are precisely \(e^{c_i}\). Tuple independence is retained,
since the test in (CSR.23) has no \(t\)-dependence. For the
\(\kappa\) uniqueness clause, on a nonzero \(E_p\) all factors
in the product of \(\lambda\)'s are unchanged by CSR11; on
the zero space every tested vector is zero. Its nonzero-test
witness is therefore also the old one.

The primitive fields \(\mathrm{PThmA}\) and \(\mathrm{PThmB}\)
are scalar operator identities on \(W_{i,p}\) and \(E_p\).
On a zero part their vectors are zero; on a nonzero part the
scalar, or the relevant product of scalars, is unchanged.
The \(\mathrm{GalD}\) field is CSR11. Positivity is preserved
by (CSR.23). All the remaining primitive fields are copied:
row/lift Hecke operators and their intertwining, commutation
and part decompositions; group unitarity/equivariance; block,
row/lift tests and \(\mathrm{liftTestB}\); the onto condition;
positive forms; both level conditions; \(\mathrm{h2P}\);
the sums, uniformity, \(\mathrm{mLat}\) and dimensions.
The unused-choice fields copy the same
\(\mathrm{ItY,ItX,ioY,ioX,Th,ioP,thF,thG}\), conjugation,
level and integral witnesses. The original \(\mathrm{Ok3}\)
and type separation are unchanged. The eleven \(\mathrm{Keep}\)
fields are the same standard character, signs, frames, beta,
principal map, group/character ties, guard and coefficient-field
value/frame/level data. This gives all of
\(\mathrm{Pins}(R',F')\), \(\mathrm{Holds}(R',F')\) and
\(\mathrm{ErrEq}(R',F')\).

Copy \(Y_0\) to a reading \(Y'\). Its fields concern the unchanged
models, levels, forms, lattices and maps. For the two central
ties, a \(\mathrm{Keep}(R',F')\) input is the same eleven-field
record on \(R_0,F_0\), so the old central tie applies. In \(Z_0\)
only the two majorant equalities inspect the changed scalar:
CSR11 proves equality of their exact \(\mathrm{XE}\) matrices.
Thus copy all of \(Z_0\) to \(Z'\), using those equalities.
The resulting thirteen identities are
\[
 \begin{gathered}
 Y'.\mathrm{tup}=Y_0.\mathrm{tup},\quad
 Y'.\mathrm{tupT}=Y_0.\mathrm{tupT},\\
 jA,\ jA0,\ eMt,\ PY,\ eFT,\ IL,\ PW,\ eE,\ \mathrm{mj}.eW,\\
 Z'.eRR,\quad Z'.jW
       \quad\text{equal their original maps}.
 \end{gathered}                                                   \tag{CSR.27}
\]
All equalities are dependent equalities on the identical original
coordinate spaces. In particular both \(\mathrm{MLat}\)-type
reading assertions and all rational tests are merely preserved
when supplied; this proves no missing original
\(\mathrm{ratG}\) or \(\mathrm{MLat}\) producer.
% END CSR12

% BEGIN CSR13
\paragraph{CSR13: the actual \(\mathrm{TwIn}/\mathrm{TwArgs}\) constructor.}
\label{claim:CSR13}
We now isolate precisely the remaining division input used by
this constructor. On the original data of CSR12, keep the
given ambient ties \(\mathrm{AmbAre}(Z_0,ex.\mathrm{amb},ex.\mathrm{amb0})\),
lift tie \(\mathrm{MjAre}(Z_0,Y_d)\) and group tie. Retain the
first output, \(\mathrm{DivDef3}\), of the source contract
\texttt{divisionCompatible\_at\_readIdentifications}, applied to
the actual updated \(R',F',Y',Z'\), the column reading below
and the specified \(sEE\). This is an exact remaining original
supplier, rather than an arbitrary \(\mathrm{TwIn}\) premise.
It includes the sharp and base row definitions and the original
theta data on these objects; its additional
\(\mathrm{DivRows3}\) conclusion is not assumed or used here.

Copy the chosen \(C_0\) of CSR9 to \(C'\) by retaining
\(\mathrm{wU,eNu,NuC}\), their bijection and lattice/metric tie.
Its point identity only sees \(D_c\), the original tuple and
\(jA\); (CSR.27) preserves all three. Thus it holds for \(C'\)
with the \emph{new} \(\mathrm{Sp},\mathrm{Zsc}\) too.
The base ambient tie is transported by the equalities of
\(\mathrm{tup},jA0\); its original ambient isomorphism is kept.
The lift tie is transported by equality of \(jW\), keeping its
actual \(eW\) and lattice certificate.

Construct the geometric half by retaining its core,
\(\kappa_{\rm geo},v_{\rm geo},\mathrm{bar},\rho_{\rm Geo},
\mathrm{equivAll},\mathrm{cTie}\) and phases. Each
Fourier--Weil tie retains its algebra equality, actual map,
coefficient-field, point, norm, character and evaluation fields.
Each original \(O_i\) and all eleven of its pin clauses concern
these same models and maps, so they are unchanged by \(\lambda'\)
and \(e'\). Likewise the seven \(\mathrm{InGeo}\) assertions
retain the same evaluation witnesses, conjugation, injectivity,
uniformity, geometric action, full lattices and dimension match.
This is transport of a supplied geometric half, not construction
of geometric periods.

Set \(x\) to have the original
\(G,\mathrm{Pl}_0,I_0,T,\mathrm{A3},\mathrm{core},\mathrm{alg},c_G\);
the \(R',F'\) of (CSR.26); this geometric half; the pins,
error and readings of CSR12; \(C'\); and the prescribed \(sEE\).
Set its division field to the retained exact \(\mathrm{DivDef3}\),
its \(JA,\mathrm{amb0}\) to the transported base ambient pair,
and its \(JM,\mathrm{mjH}\) to the transported lift pair.
These are all the fields of \(\mathrm{TwIn}\).
Define \(a:\mathrm{TwArgs}\) with these same
\(D_c,\mathrm{Sp},\mathrm{Zsc},\lambda',e',x\).
Its whole-package equality (CSR.14) is literal;
\(\mathrm{lam\_eq}\) is CSR11 on \(E_p\ne0\);
\(\mathrm{intersectionNumber\_baseChange}\) is the copied
\(\mathrm{Holds}\). Its geometric \(\kappa,v_0,\mathrm{fw}\)
equalities are those just retained, and every remaining map
equality is (CSR.27).
Thus the original constructor has
\[
 \begin{gathered}
 a.D_c=D_c,\quad a.\mathrm{Sp}=\mathrm{Sp},\quad
 a.\mathrm{Zsc}=\mathrm{Zsc},\\
 a.x.sEE=sEE,\quad
 a.x.C.\mathrm{wU}=C_0.\mathrm{wU},\quad
 a.x.C.\mathrm{eNu}=C_0.\mathrm{eNu},
 \end{gathered}
 \tag{CSR.28}
\]
and retains \(\mathrm{InGeo}\) on the same geometric half.
The statement is conditional exactly on the supplied original
readings/geometric half, (CSR.22) and this original
\(\mathrm{DivDef3}\) supplier.
% END CSR13

% BEGIN CSR14
\paragraph{CSR14: the original selected realization with one period update.}
\label{claim:CSR14}
Fix the original \(q,R,\mathrm{RecPinE}(q,R)\). Transport the
whole \(q.\mathrm{sat}.P,q.\mathrm{sat}.S\) along the column
equality, obtaining \(sP,sS\), and set \(S=sS.\mathrm{toSTiesR}(R)\).
Its \(\mathrm{fin.inHH}\) supplies every \(pw_u=\mathrm{pwOf}(S,u)\).
Take the face and \emph{one} \(C_0\) of CSR9; put
\(w=\!C_0.\mathrm{wU}\) and \(e_u=\!C_0.\mathrm{eNu}_u\).
Their original meanings are (CSR.12).
At \(Z_s=zsOf(R_0,F_0)\), define the period object
\[
 \mathrm{Sp}^{*}=\mathrm{spOf}
 \left(S.X,s,D_c,S.Y.hK,S.G,w^{-1},
            (pw_{w^{-1}(m)})_m,(b_m)_m\right),
 \tag{CSR.29}
\]
where \(s=\mathrm{twAt}(D.\mathrm{embLL},\mathrm{A3}.v_0)\) and
\[
 b_m(i)=D.eMu_{w^{-1}(m)}
        \bigl(D.jN_{w^{-1}(m)}
          (e_{w^{-1}(m)}(e_i))\bigr).
 \tag{CSR.30}
\]
The index \(i\) in (CSR.30) is transported by
\(w(w^{-1}(m))=m\); this is the exact basis used in
\(\mathrm{spRead}\). Replace \(\mathrm{Sp}_0\) only at this
one \(Z_s\) by \(\mathrm{Sp}^{*}\), leaving its other values
as in (CSR.15).

For the supplied natural pair \((\Delta,h_0)\) and any real \(c'\),
use the original scale function
\[
 \begin{gathered}
 z=\mathrm{ZscOfNum3}
       (\mathrm{ofRecord}(D,q.O.L,\Delta,h_0),
        \mathrm{id},q.O.L.N_c),\\
 z'.l6m(j)=z.l6m(j)-c' z.l6u,\quad z'.l6u=0 .
 \end{gathered}
 \tag{CSR.31}
\]
This is the literal \(\mathrm{Fc.zsc}\), evaluated on each
subspace. The original \(q.\mathrm{sat}\) supplies
\(\mathrm{EvalTieT}\) at its own \(t,c\), with \(c>0\);
\(R_0.\mathrm{models\_HE}\) is its retained model certificate.
CSR10 supplies \(s_E,\kappa>0\) and (CSR.21).
By the definition of (CSR.29), its Hilbert space is exactly
the target of \(s_E\) after the stated twist, so \(sEE=s_E\)
by the resulting type equality, with no second linear map.

Apply CSR12--CSR13 with these \(\mathrm{Sp},z',sEE\).
Their actual upstream support and division suppliers remain
as stated there. Now apply CSR9 to this \(a,x=a.x\);
the two \(\mathrm{ReadC}\) maps in (CSR.28) are its required
equalities. We obtain all five column conclusions (CSR.17).
The whole-package equality carries (CSR.31) to
\[
        x.F.\mathrm{Zsc}
             =\mathrm{Fc.zsc}(\mathrm{divisionPair}(x),
                                      \Delta,h_0,c').
 \tag{CSR.32}
\]
It also carries (CSR.29) to the exact period object in
\(\mathrm{RealIn.spPin}\): both the inverse word map and every
vector \(b_m(i)\) agree by (CSR.28). This explains why equality
of \(eNu\), although unused in CSR9's last step, is necessary here.
Finally transport \(s_E\) along the same package equality.
Its formula and the factor \(n_P\) in (CSR.21) retain
\(q.\mathrm{sat}.t,q.\mathrm{sat}.c\), the same good operator,
and the original level subspace of \(\mathrm{MMap}\).
This proves the full selected-realization packet, with the
remaining suppliers precisely those of CSR13.
% END CSR14

% BEGIN CSR15
\paragraph{CSR15: the standard character, point and rational exponent transports.}
\label{claim:CSR15}
Two continuous adelic additive characters, trivial on the
diagonal field, agree if they agree on the archimedean factor.
Their values on a rational element in the finite factor are the
inverses of their common archimedean value. Integral finite
adeles are limits of integral rational elements by the Chinese
remainder theorem. Equality on them follows by continuity.
For an arbitrary finite adele, choose a diagonal field element
with the same finitely many principal parts; the difference is
integral at every finite place. Such an element is obtained by
clearing the finite list of denominator ideals and applying
the Chinese remainder theorem to that fractional ideal.
Additivity gives equality on the finite factor, and then on
the whole adele group.

In \(\mathrm{IsStdChar}(D,\psi,-1)\) the archimedean formula is
the same specified trace exponential. Thus the character in
\(x.\mathrm{pins.keep.std}\) equals the one in
\(R.T.\mathrm{theta}.W\), transported through
\(\mathrm{RecPinE}\); this is the literal \(\mathrm{psiPin}\).
No class-field-theoretic uniqueness assertion is needed here.

The original spectral statement consists of \(\mathrm{SeeAt}\),
\(\mathrm{EvalTieT}\) and \(\mathrm{RowJ}\). It uses the
representation, level, evaluation point, scalar, measure and
the actual \(eE\) reader. None of these changes under the
\(\lambda,e'\) update. The whole-package equality and
\(a.eE\) therefore transport it to \(x\) with the same
\(q.\mathrm{sat}.t,q.\mathrm{sat}.c\).
For the rational exponent data, set
\[
 \begin{gathered}
 S_x=q.\mathrm{core.row}.\mathrm{ratDatR},\quad
 n=q.\mathrm{core.row.nexp},\\
 n'=(n.fi,n.geq,n.rho,n.chi,n.eF,n.bsup,n.act).
 \end{gathered}
 \tag{CSR.33}
\]
These seven fields do not use the updated scalar or error.
Consequently they define the actual
\(\mathrm{NExpRealS}(x.R,x.c_G,S_x)\), with exactly the original
\(eF\) and \(fi\).
Finally the adelic isomorphism
\(\iota_S=\mathrm{iotaX}(sP.X,sP.\mathrm{inst})\)
is the fixed rational-embedding map followed by the group cast
for \(sP.\mathrm{inst}\). After that equality of instances,
the cast is the identity; it is therefore exactly the
\(\mathrm{HEq}\) required by \(\mathrm{RealIn}\).
% END CSR15

% BEGIN CSR16
\paragraph{CSR16: every field of the original \(\mathrm{RealIn}\) and \(\mathrm{RealTB}\).}
\label{claim:CSR16}
Retain the original host certificate, positive-definite \(V_i\),
and the standard frame/lattice equalities for \(V_i\) and \(X'\).
These are the corresponding fields of the source's
\(\mathrm{hostHyp\_frames}\) supplier at \((q.w,q.Z)\);
their uniform production is not inferred from a record cast.
On the selected packet of CSR14 define \(y:\mathrm{RealIn}(x,R,S)\)
as follows. Use CSR15 for \(\mathrm{psiPin}\), (CSR.33) for
\(S_x,\mathrm{nexp}\), and the retained host data for
\(\mathrm{stdV,stdX,host,vpos}\).
Take \(\mathrm{hin}\) from the same \(sP.\mathrm{fin}\) and
\(\mathrm{spPin}\) from (CSR.29)--(CSR.30). Set
\[
 t_S=q.\mathrm{sat}.t,\quad c_S=q.\mathrm{sat}.c,\quad
 \kappa_S=\kappa,\quad
 \iota_S=\mathrm{iotaX}(sP.X,sP.\mathrm{inst}),\quad
 u_S=1,\quad sEE_c=s_E .
 \tag{CSR.34}
\]
The \(\mathrm{instEq}\) and \(\mathrm{iotaS\_eq}\) fields are
CSR15, and \(sEEj\) is the exact map equality in CSR14.
The \(\mathrm{seeLev}\) implication may ignore its existential
antecedent because the transported \(\mathrm{SeeAt}\) already
provides the specified conjugating point \(t_S\). The
\(\mathrm{seeEval}\) implication likewise uses the already
supplied evaluation scalar \(c_S\). The norm of \(u_S=1\) is one.
The \(\mathrm{seeDef}\) field is (CSR.21), and
\(\mathrm{seeK}\) consists of \(\kappa>0\) and its exact
\(n_P\)-norm identity. That identity is already proved, so its
additional existential normalization antecedent is unnecessary.
This checks every field of \(\mathrm{RealIn}\).

The original \(\mathrm{realTOf}(y)\) retains \(x.R,x.sEE\),
its characteristic, trace and standard-character pins, its two
subspaces and column/period data, its frames and all-place
conditions. Its word inverse is \(x.\mathrm{wd}=(x.C.\mathrm{wU})^{-1}\);
the two inverse identities and its \(bR\) vectors are exactly
(CSR.30). All remaining real fields are the above fields of \(y\).
The original \(\mathrm{realTBOf}(y)\) additionally retains
\(x.\mathrm{core},x.\mathrm{geo}.\kappa,v_0,\mathrm{fw},
\mathrm{bar},O,\mathrm{pins},\rho_{\rm Geo},\mathrm{equivAll},
\mathrm{cTie}\). Its phase is the original phase because
\(x.\mathrm{pins.phs}(i)=1\) gives
\(\mathrm{phEff}(x.R)(i)=1\) at every \(i\).
Thus its six output equalities are literal:
the same \(t_S,c_S,S_x,eF,fi\), and \(u_S=1\), with the
unchanged geometric half. In particular no new nonzero geometric
witness, rational form or analytic representative is hidden in
this transport.
% END CSR16

% BEGIN CSR17
\paragraph{CSR17: the terminal-coordinate factor cancels in the actual lower map.}
\label{claim:CSR17}
Put \(N=n+1\), let \(\lambda\) be a dominant weight of length
\(N\), and let \(\gamma\) interlace it. In the source's notation,
\[
 b_i=\lambda_i-\gamma_i\geq0\quad(1\leq i\leq n),\qquad
 n_{\rm Br}=\sum_{i=1}^{N}\lambda_i-\sum_{i=1}^{n}\gamma_i
       =\lambda_N+\sum_{i=1}^{n}b_i .
 \tag{CSR.35}
\]
For each embedding \(\tau\) take a nonzero scalar \(a_\tau\),
and put \(A_\tau=\operatorname{diag}(1,\ldots,1,a_\tau)\).
The polynomial action includes its determinant factor:
\(\operatorname{actP}_{\lambda}(A)\) multiplies a monomial by
\[
 \prod_\tau a_\tau^{\lambda_{\tau,N}
                 +\sum_j \alpha_{\tau,Nj}},
 \tag{CSR.36}
\]
where \(\alpha\) is its exponent array. The operator
\(\mathrm{coefLast}\) in the branching definition keeps exactly
the monomials with last-row exponents \(b_{\tau,j}\) for
\(j\leq n\) and zero last-column exponents; all other monomials
are discarded. On every kept monomial (CSR.36) is
\(\prod_\tau a_\tau^{n_{\rm Br}(\lambda_\tau,\gamma_\tau)}\).
It is therefore a common scalar before the remaining linear
projection onto the prescribed coefficient space and before
the fixed branching normalization. Monomial linearity proves
the exact identity
\[
 \left(\prod_\tau
 a_\tau^{-n_{\rm Br}(\lambda_\tau,\gamma_\tau)}\right)
 \operatorname{brStdP}_{\lambda,\gamma}
       (\operatorname{actP}_{\lambda}(A)P)
   =\operatorname{brStdP}_{\lambda,\gamma}(P)
 \tag{CSR.37}
\]
for every polynomial \(P\). The conversion of \(b_i\) to a
natural exponent uses \(b_i\geq0\); no sign or determinant
shift is dropped.

Now use the original terminal column. Write
\(\pi=X.\pi_T(u)\), so \(\pi\bar\pi=T_u\), and
\(z_0=\pi e_N\). The standard coordinates from TG12 give
\[
 b_Z=\operatorname{diag}(1,\ldots,1,\bar\pi^{-1}),\qquad
 m_Z(\tau)=\operatorname{diag}
      (1,\ldots,1,\operatorname{emb}_{\tau}(\bar\pi)).
 \tag{CSR.38}
\]
Indeed the first \(N-1\) columns of \(b_Z\) are the standard
inclusion and its last column is \(z_0/T_u\). The relation
\(\pi\bar\pi=T_u\neq0\) proves both (CSR.38) and invertibility.
Set \(a_\tau=\operatorname{emb}_{\tau}(\bar\pi)\).
The transpose of \(A_\tau\) equals \(A_\tau\), so the
definition of \(\mathrm{fnOf}\), including its determinant
shift, gives
\(\mathrm{fnOf}(m_Z,P)=
\mathrm{fnOf}(1,\operatorname{actP}_{\lambda}(A)P)\).
Applying \(\mathrm{polOf}\) in the original standard frame
recovers this exact polynomial.

The definition of the original lower map is
\(\mathrm{brAt}(a)\); its prefactor is precisely that on the
left of (CSR.37). Hence its value is
\(\mathrm{fnOf}(1,\mathrm{brTen}(D,u)P)\), with the original
target weight and branching constant. This is
\(\mathrm{LowerMapTieS}\). For a nonbase used tuple, the
actual twisted comparison TX14--TX16 identifies the same
evaluations, weights and frame arguments with their standard
ones. Substitution in (CSR.37) therefore gives
\(\mathrm{LowerMapTieTw}\) for every tuple in its stated
domain, including its \(1\leq m\) guard. The full
\(\mathrm{RecPinE}\) transports these statements to
\(R.T.\mathrm{cols}\) and its original twisted forms.
This proof uses the already reconstructed standard and
twisted coordinate identifications; it adds no branching
existence or arbitrary intertwiner.
% END CSR17

% BEGIN CSR18
\paragraph{CSR18: the complete row-side determination record.}
\label{claim:CSR18}
For any original finite lattice class, forgetting its
auxiliary residue matrices maps its stabilizer into the
integral automorphism group of its underlying lattice \(N\).
This map is injective. If its rational automorphism is the
identity, the auxiliary matrix \(M\) fixes every vector of
the chosen residue basis: \(M b(v)=b(v)\). Surjectivity of
that basis map gives \(M=1\), at every auxiliary place.
Thus the entire stabilizer element is the identity.
Lagrange's theorem yields
\[
 |\operatorname{Stab}(c)|\mid |\operatorname{Aut}_{\mathcal O}(N)|
       \mid D_{\Gamma}.
 \tag{CSR.39}
\]
The final divisibility is the actual finite least common
multiple of TL2. The ranks of the row lattices and of every
lift lattice occur among its specified ranks; the source's
\((D.\mathrm{mem}\,D.u_0).C.DGam\) is this same integer.
This proves both universally quantified divisibilities in
the \(\mathrm{gamDvd}\) field, for every tuple and
every class, not only the base tuple.

Construct \(\mathrm{RowSideDetermined}(R.T,P)\), where
\(P\) is the original division pair of \(x\).
Its first two fields are CSR17. Its \(\mathrm{buDef}\) field
is \(x.Y.\mathrm{bu}\). Its sharp ambient field is the original
ambient identification in \(q.x.\mathrm{obj.amb}.1\),
transported through the whole packet of CSR8 and the unchanged
tuple and \(j_A\) maps of CSR12--CSR13. This is the sharp
ambient \(\mathrm{ex.amb}\); replacing it by
\(\mathrm{ex.amb0}\) would not typecheck the field.
Its \(\mathrm{faces}\) witnesses are the full functions
\(x.F.Dc,x.F.Sp\), whose evaluations at \(P.Zs\) are
definitionally \(P.Dc,P.Sp\). Finally use (CSR.39) for
\(\mathrm{gamDvd}\). All six fields are now supplied.
The argument transports the given ambient and analytic
data; it does not manufacture the original \(\mathrm{MMap}\)
or a geometric ambient witness.
% END CSR18

% BEGIN CSR19
\paragraph{CSR19: finite eigensystem selection and the exact \(\mathrm{RecIn}\).}
\label{claim:CSR19}
There are two elementary reading identities. Apply the
original \(x.Y.\mathrm{mtRead}\) at a used tuple \(\mu\),
put \(t=x.Y.\mathrm{tupT}(\mu)\), and substitute the inverse
of \(x.Y.eMt\). The two tuple inverse identities and
cancellation of this linear equivalence give the required
\(\mathrm{MtRead}\) for the original \(\mathrm{twPOf}\);
its used-tuple branch is the branch selected by \(\mu\).
For \(\mathrm{MjRead}\), use \(x.Y.\mathrm{mjRead}\) on
the vector \(x.Z.jW(i,v)\). The image identity for
\(x.Y.\mathrm{mj.eW}\) places it in the original good
subspace, where the good projection is the identity.
Substituting \(x.Z.\mathrm{jW\_read}\) gives exactly
\(\mathrm{jDT}(\mathrm{realTOf}(y))\). These substitutions
retain the original tensor map and multiplicity map.

For the eigensystem part retain the following exact
producer outputs on this same \(x.R,x.F,x.Y\). Every
original type \(\pi\) has its spherical eigensystem
\(\alpha_\pi=\mathrm{eigOfType}(R,F,Y,\pi)\).
Every such eigensystem has a carrier \(p\).
A carrier contains the corresponding nonzero isotropic
part under the original \(\mathrm{mj.eW}\), and gives
the original row-part and lift-part containments.
These are the outputs of
\(\mathrm{exists\_sphericalEigensystem\_firstLift\_baseLevel}\),
\(\mathrm{carries\_of\_hasEig}\),
\(\mathrm{isot\_le\_part}\),
\(\mathrm{prt\_le\_hPart\_of\_carries}\), and
\(\mathrm{prtW\_le\_hPart}\), respectively.
Their inputs include the original host and nonzero
character-level idempotents. Also retain the all-embedding
identity
\[
 \begin{split}
 &D.\mathrm{isot}(i+1,\pi)\ne0,\quad
 D.\mathrm{isot}(i+1,\pi)
       \subseteq (x.R.Wp3(i,p))\circ(x.Y.\mathrm{mj.eW}(i))
 \\
 &\hspace{12mm}\Longrightarrow\
 q.O.X.\mathrm{logLam}(i+1,\pi,e)
   =\log x.R.\mathrm{lam3}(i,\mu_0,p,\mathrm{InfinitePlace}(e)).
 \end{split}
 \tag{CSR.40}
\]
The subspace on the first line is a comap. The
off-base-embedding reciprocity assertion in (CSR.40) is
an explicit retained input; it is not implied by finite
eigensystem linear algebra.

Here is the complete finite construction from these
inputs. Let \(\Pi\) be the image of \(\pi\mapsto\alpha_\pi\)
on the original finite set of types. Choose one carrier
\(p_\alpha\) for each \(\alpha\in\Pi\), and set
\[
 cOf(\pi)=\alpha_\pi,\qquad
 piOf(\alpha)=\alpha,\qquad
 labPi(\alpha)=\mathrm{toIx2}(p_\alpha).
 \tag{CSR.41}
\]
The first map lands in \(\Pi\); the second is injective.
If \(cOf(\pi)=cOf(\pi')\), the common eigensystem has
one common carrier. Apply (CSR.40) to both nonzero
isotropic parts at that carrier. Their logarithmic
parameters agree at every embedding. This proves
\(\mathrm{EigIdx.lam\_const}\), with its original
nonzero guards and range \(1\leq j\leq m\).
The two retained part containments for \(p_\alpha\)
give the remaining two fields of \(\mathrm{PartsRec}\);
the other two are the definition of \(cOf\) and
injectivity of \(piOf\). No carrier is reselected in
the part comparisons.

It remains to install the character and tuple data.
Unpacking the whole \(\mathrm{TwArgs}\) packet and
\(\mathrm{RecPinE.s.chr}\) gives
\(x.R.C=R.T.n.cf\). Transport \(R.T.n.tp\) across exactly
this equality to define \(Tp\). Transporting a dependent
record along its base equality does not alter its
\(\xi\)-function. The original \(\mathrm{twArgs\_lev}\)
identity and \(q.x.n.\xi\) therefore give, for every
used tuple,
\[
 x.R.\xi_H(x.Y.\mathrm{tupT}(\mu))=Tp.\xi(\mu).
 \tag{CSR.42}
\]
The eight fields of \(\mathrm{RecIn}\) are now
\(E,labPi,piOf,Tp,tpEq,tupXi,mtRead,mjRead\), supplied
by (CSR.41), this transport, (CSR.42), and the first
paragraph. The \(\mathrm{realRecOf}\) constructor copies
the original characteristic, tied readings, three part
identities, character, tuple inverses, finite isometries
and frame matching. Its \(\mathrm{PartsRow}\) is precisely
the four-field \(\mathrm{PartsRec}\) just constructed.
Thus the finite assembly is proved while its
eigensystem and all-embedding parameter producers remain
explicit.
% END CSR19

% BEGIN CSR20
\paragraph{CSR20: selected primes and both distinguished contractions.}
\label{claim:CSR20}
Let \(HD\) be an original host datum on
\(R.T.\mathrm{toTiedDataS}\), of type \(\mathrm{HostDatumC}\), with
\(HD.d3b=q.x.\mathrm{pin}.d3b\).
The standard-kernel seam in \(\mathrm{RecPinE}\) is the
identity on the original field. Consequently its prime
map is the identity, and the image of
\(\mathrm{QE}(R.T.\mathrm{ker.sel},R.T.\mathrm{ker.tup}(\mu))\)
under that map is exactly \(I.Q(\mu)\).
This follows on ideals by inverse image under the identity;
the finite-set equality follows element by element.

The original \(\mathrm{Pq\_over}\) gives the first
over-prime condition. Any prime in the above image is
one of the \(I.\mathrm{sel}(v)\), by the definition of
\(I.Q\). Since \(q.\mathrm{core.sel}\) identifies the
ranges of \(I.\mathrm{sel}\) and \(q.\mathrm{core.a.Pq}\),
the selected-prime range condition follows. If a prime
is presented through the source's carried-prime relation,
its ideal equals the prime-map ideal, so equality of
prime ideals gives the same conclusion.

For a used tuple, the original tuple inverse
\(\mathrm{ixOf}(Y_0.\mathrm{tupT}(\mu))=\mu\),
the old level equality and the unchanged \(\mathrm{tupT}\)
give
\[
 x.R.Q(x.Y.\mathrm{tupT}(\mu))=I.Q(\mu).
 \tag{CSR.43}
\]
This is the second field of \(\mathrm{SelSideRT}\).
For the third, suppose \(p_{\mathrm{de}}\) lies in
the ideal \(q.\mathrm{core.a.Pq}(v)\). Write that
selected prime as \(I.\mathrm{sel}(v')\).
Its contraction to \(F\) is \(v'\), by the original
\(\mathrm{above}\) pin, so \(v'\) lies over
\(p_{\mathrm{de}}\). The original \(d3b.P_{\rm mem}\)
and \(\mathrm{selectedPlacesAgree}.1\) now identify its
contraction to \(K\) with \(d3b.pK\). Replacing
\(d3b\) by the equal \(HD.d3b\) proves the required field.
The fourth field follows by the same argument with
\(p'_{\mathrm{de}},P'_{\rm mem}\),
\(\mathrm{selectedPlacesAgree}.2\), and \(pK'\);
both distinguished primes are checked separately.
This completes all four fields of \(\mathrm{SelSideRT}\)
for the same \(x,y,r,HD\). It is conditional on the given
host datum and does not assert its class-field-theoretic
existence.
% END CSR20

% BEGIN CSR21
\paragraph{CSR21: the original seven-field bundle and every \(\mathrm{BWit}\) field.}
\label{claim:CSR21}
Fix \(q,R,\mathrm{RecPinE}(q,R),c\) and the source's
integer pair \(DH=(\Delta,h_0)\) with its exact equalities
\[
 D.s.Del=\Delta,\qquad q.O.L.\mathrm{etaSz}=D.R.m\,h_0.
 \tag{CSR.44}
\]
Use a single selected packet from CSR14, a single
\(\mathrm{RealIn}\) from CSR16, a single \(\mathrm{RecIn}\)
from CSR19 and the same host datum as in CSR20.
Put \(P=\mathrm{divisionPairOfAdelicFormsFaceColumnsData}(x)\).
The three equalities from the selected construction are
\[
 C.c.\mathrm{face}=x.F.Dc,\qquad
 C.Sp=x.F.Sp,\qquad
 \mathrm{zsc}(P,\Delta,h_0,c)=x.F.Zsc.
 \tag{CSR.45}
\]
The middle equality is dependent equality, namely the
source's \(\mathrm{HEq}\) over the first. Transport
\(\mathrm{realTBOf}(y)\) through these three equalities
together, obtaining \(B'\). Define the literal bundle
\[
 b=\bigl(R,\ sS.\mathrm{toSTiesR}(R),\
             P,\ \mathrm{CSR18},\ DH,\ C,\ B'\bigr).
 \tag{CSR.46}
\]
These are all seven fields of \(\mathrm{Bundle}\);
in particular \(b.R=R\) is definitional.

The complete package \(RPk\) consists of
\(P,Dc,Sp,Zsc,B\). After (CSR.45) and the corresponding
transport of \(B\), its two realizations \(\mathrm{rpkB}(b)\)
and \(\mathrm{rpkX}(y)\) are equal. One can check this by
successively replacing equal bases; when all three
bases agree the last transport is the identity and
all five fields agree. This proves equality of the
whole package, not just a list of isomorphic vector
spaces. It supplies the equality field of
\(\mathrm{BTw}\); the other fields are the same
\(\mathrm{Holds}\) from \(a\), \(y\), \(r\), and \(HD\).

The remaining witness is now explicit. Its fields
\(x,a,ax\) are the selected tuple and its actual
\(\mathrm{TwArgs}\) equality. Its \(y\) is this
\(\mathrm{BTw}\), and its \(\mathrm{recPin}\) is the
initial \(\mathrm{RecPinE}(q,R)\).
Its \(s\) field uses the same transported \(sP,sS\),
with both equalities to \(q.\mathrm{sat}.P,S\);
the equality \(b.S=sS.\mathrm{toSTiesR}(b.R)\)
is definitional. Its \(pt\) field is
\((y.t_S=q.\mathrm{sat}.t,\ y.c_S=q.\mathrm{sat}.c)\).
Its \(sx,eF,fi,uS\) fields are exactly the six
equalities checked in CSR16. Its \(\mathrm{parts}\)
is CSR19. Its \(cW\) is the pair of actual
weight and \(X_g\) equalities from CSR14.
Its \(d3b\) is the given equality for \(HD\);
its \(\mathrm{side}\) is CSR20 and its
\(\mathrm{colPins}\) is CSR7 transported by CSR14.
This is every field of the original \(\mathrm{BWit}(q,b)\).
The construction neither changes \(R\) nor picks
separate bundles for its different requirements.
% END CSR21

% BEGIN CSR22
\paragraph{CSR22: the eventual same-record \(\mathrm{BunU}\) implication.}
\label{claim:CSR22}
Fix the original \(h_L,N,\Gamma\) and
\(p=\mathrm{prmR}(N,\mathrm{SExt.none},\Gamma)\).
Keep the universally quantified original extension
premise. Write
\(\mathfrak a=(K,B,f,T,w,Z,Wt,p_0,t,nch,n,xd,c)\)
for its full dependent tuple of arguments. Then the premise is
\[
 \forall\mathfrak a,\qquad
 N(\mathfrak a)\ \Longrightarrow\
 \mathrm{Nonempty}\bigl(\mathrm{NEb}(h_L,\mathfrak a)\bigr).
 \tag{CSR.47}
\]
It is retained wherever the original upstream producer
uses it; separate output-specific extensions are not
substituted for (CSR.47).

The sufficient producer interface is now concrete.
At each original \(q,R,\mathrm{RecPinE}(q,R)\), retain
the host/frame packet used in CSR16, the nonzero-part
support condition (CSR.22) used in CSR11, the actual
first division-compatibility output used in CSR13,
the five eigensystem outputs and all-embedding identity
of CSR19, and a host datum with the exact \(d3b\) pin
of CSR20. The support condition is the literal output
needed from the original multiplicity equality; no
new multiplicity assumption on a different realization
is sufficient. The original source obtains these
packets from its named analytic, geometric, multiplicity,
reciprocity and extension hypotheses. Their names alone
are not proofs of their production here.

The conditional analytic construction SPR25 can supply
the downstream original multiplicity map on the same
record once its stated initial bounded extremal seed
and target \(\mathrm{FinPair}\) inputs, and its other
original guards, have been supplied. This reduces the
analytic producer interface; it does not erase those
initial inputs or create the common geometric record.
All occurrences of \(\mathrm{MMap}\) in CSR1--CSR21
refer to that original map.

For precision, the source's eventual quantifier has
the following order:
\[
\begin{split}
 &\exists\varepsilon_0>0\
 \forall\varepsilon,C_8,C_5,c_6,C_b,\
 [\,0<\varepsilon\leq\varepsilon_0,\ 
 \varepsilon<1/52,\ 1\leq C_8\,]\\
 &\quad\exists K_0\ \forall C,\
 [\,\mathrm{Std}(C),\ C.P.Cpde=C_8,\
 C.\mathrm{Large}(K_0,C.P.CF)\,]\\
 &\qquad\exists c_1\ \forall c_D\geq c_1\
 \exists c_2\ \forall c_h\geq c_2,\
 \left[\kappa.\mathrm{Pos}\Rightarrow
 \exists x_0\ \forall K,\
 \kappa.\mathrm{Of}(K)\Rightarrow
 x_0\leq K.s.x\Rightarrow\cdots\right],
\end{split}
\tag{CSR.48}
\]
where \(\kappa=(\varepsilon,C,c_D,c_h,C_5,c_6,C_b)\)
and \(\kappa.\mathrm{Pos}\) means that its last three
constants are positive. Assume the retained producer
packets hold in this order. In the original theorem
the two eventual inputs are the host/frame packet
and \(\mathrm{MultEq}\); the remaining suppliers are
pointwise under the displayed original premises.
Our conditional interface also permits their exact
outputs to be supplied eventually in (CSR.48).

Finite conjunction preserves (CSR.48). Take the minimum
of the positive upper bounds for \(\varepsilon\), the
componentwise maximum \(K_0\) of the size packages,
and the maximum of the \(c_1\)'s. After \(c_D\) has
been fixed, take the maximum of the resulting \(c_2\)'s;
after \(c_h\) and \(\kappa.\mathrm{Pos}\) are fixed,
take the maximum of the resulting \(x_0\)'s.
The original \(\mathrm{Large}\) inequalities are
monotone under the componentwise maximum. Each
subsequent universal variable therefore satisfies
every required guard. No point, tuple, word or
embedding is selected before its uniform threshold.

At every resulting late \(K\), for every
\(q,R,\mathrm{RecPinE}(q,R),c,DH\), apply CSR21.
It returns
\[
 \exists b:\mathrm{Bundle}(q.w.D\,q.Z,q.O,c),\qquad
 b.R=R\quad\hbox{and}\quad
 \mathrm{Nonempty}(\mathrm{BWit}(q,b)).
 \tag{CSR.49}
\]
The additional original guards \(q.O.L.sc=K.s\) and
\(\mathrm{Val}(K.C,K.cD,K.ch,q.O.L)\) are retained;
the finite final join does not need to use them anew.
Thus (CSR.49), with exactly (CSR.48), is the literal
\(\mathrm{BunU}(p)\) conclusion through
\(\mathrm{SelAll.bwit\_exists}\), conditional on the
identified original producers.

Crucially, CSR9 selected one \(C_0:\mathrm{ReadC}\)
before any later \(T,a,x\); every later constructor
retains its \(w_U,e_\nu\). If an actual
\(\mathrm{RawSel}\) is subsequently chosen, its six
bundle projections must all come from that one
selector at \((\mathcal F,x,i)\), with \(x>\mathcal F.x_1\).
This theorem does not choose independent projections,
prove existence of the common family, construct
\(\mathrm{MemberE}\), or prove the additional
\(\mathrm{ExBL}\) and \(\mathrm{RowL}\) producers.
% END CSR22

\paragraph{Exact remaining boundary.}
The finite column constructor, normalization, scalar
update, geometric copying, lower-map pin, selected-prime
comparison, finite eigensystem assembly and complete
same-record bundle join are now proved with their
literal inputs. The common selected point and its
geometric/readings data, the initial analytic inputs,
multiplicity support, original division installation,
spherical/carrier/part production, all-embedding
reciprocity, and host-datum production remain where
explicitly retained above. In particular, preserving
the old \(\mathrm{mLat}\) field under a scalar update
does not establish the separate \(\mathrm{ratG/MLat}\)
branch, and the first division output does not supply
all of \(\mathrm{DivRows3G}\).
